\chapter{Results and Discussion}
\thispagestyle{empty}

All results in this chapter are obtained on people not seen during training, with the protocol of
Chapter~4: leave-one-subject-out for DS-1 (14 folds) and stratified grouped 10-fold
cross-validation for DS-2, with nested tuning on the training subjects. Balanced accuracy is per
window; ``subjects'' or ``children'' counts the people
classified correctly by majority vote.

\section{Main results}
Table~\ref{tab:main} gives the results of the proposed pipeline (SMVMD common modes, Riemannian
tangent space and LightGBM, with the mode set chosen by the inner cross-validation).

\begin{table}[H]
\caption{The proposed pipeline on unseen people. Window-level balanced accuracy with 95\% CI over
subjects, AUC, sensitivity and specificity (\%), and people classified correctly by majority
vote.}
\label{tab:main}\centering
\gentable{tab_main}\\[2pt]
{\footnotesize \dstwonote}
\end{table}

\textbf{ADHD (DS-2).} The pipeline reaches \R{MainDStwoBal}\% balanced accuracy (95\% CI
\R{MainDStwoBalCI}\%) and an AUC of \R{MainDStwoAUC} on the windows of unseen children, and
classifies \R{MainDStwoSubj} children correctly (\R{MainDStwoSubjAcc}\%). Sensitivity (\R{MainDStwoSens}\%) is
higher than specificity (\R{MainDStwoSpec}\%): control windows are misclassified more often than
ADHD windows. In every outer fold the inner cross-validation chose all four common modes
together (\R{MainDStwoChosen}).

\textbf{IDD (DS-1).} All 14 participants are classified correctly in all three settings, with
balanced window accuracies of \R{MainDSoneRestBal}\% (rest), \R{MainDSoneMusicBal}\% (music) and
\R{MainDSoneRestMusicBal}\% (rest and music). Here the inner
cross-validation mostly chose the $\alpha$ mode alone (Table~\ref{tab:chosen}): the two
disorders are detected from different frequency ranges by the same pipeline. Because the two DS-1
groups are also separated by a single property of the rest recordings, their relative beta power
(AUC \R{DSoneRestRelBetaAUC}, Section~\ref{sec:datachecks}), and have only seven adults each,
the DS-1 result shows that the pipeline transfers to a second disorder, not that it has found an
EEG marker of IDD.

\begin{table}[H]
\caption{Common-mode set chosen by the inner cross-validation in the outer folds.}
\label{tab:chosen}\centering
\gentable{tab_chosen}
\end{table}

\section{Choice of classifier}
Table~\ref{tab:classifiers} and Fig.~\ref{fig:classifiers} compare four classifiers on the same
tangent-space features of DS-2, for the SMVMD common modes and, for reference, for the broadband
covariance of the undecomposed window (the representation of Step~7).

\begin{table}[H]
\caption{DS-2: classifier comparison on tangent-space features. Balanced window accuracy (\%),
AUC, and children correct.}
\label{tab:classifiers}\centering
\gentable{tab_classifiers}\\[2pt]
{\footnotesize \dstwonote}
\end{table}

\begin{figure}[H]\centering
\genfig[0.9\textwidth]{classifier_comparison}
\caption{DS-2: balanced accuracy of the four classifiers on three tangent-space representations.
Error bars: 95\% CI over children; dashed line: chance.}
\label{fig:classifiers}
\end{figure}

Gradient boosting improves on the linear and kernel classifiers for the SMVMD common modes: with
the mode set chosen in the training folds, logistic regression reaches \R{AblSelectLrBal}\%, the
SVM \R{AblSelectSvmBal}\%, XGBoost \R{AblSelectXgbBal}\% and LightGBM \R{AblSelectLgbmBal}\%, and
LightGBM classifies the most children correctly (\R{AblSelectLgbmSubj}). On the broadband
covariance, LightGBM also gives the highest score (\R{AblRawLgbmBal}\%, \R{AblRawLgbmSubj}
children), while XGBoost is there slightly below the SVM. All these differences are small, a few
points at most; the classifiers should be regarded as close to equivalent, with LightGBM the
highest on both representations.

\section{Contribution of the modes and of the decomposition}
Table~\ref{tab:modes} and Fig.~\ref{fig:modes} show, for DS-2, the result of each common mode on
its own, of all four, and of the sum of all modes, and compare the SMVMD common modes with
zero-phase Butterworth band-pass filters with the same band edges, under an otherwise identical
pipeline (LR and SVM, which were run for every variant).

\begin{table}[H]
\caption{DS-2: balanced window accuracy (\%) of single and combined components, SMVMD common
modes against Butterworth bands with the same edges.}
\label{tab:modes}\centering
\gentable{tab_modes}\\[2pt]
{\footnotesize \dstwonote}
\end{table}

\begin{figure}[H]\centering
\genfig[0.85\textwidth]{mode_comparison}
\caption{DS-2, SVM: balanced accuracy of each SMVMD common mode and Butterworth band, and of all
four together. Grey line: undecomposed broadband covariance (Step~7). Error bars: 95\% CI.}
\label{fig:modes}
\end{figure}

Three observations follow. First, the slow modes carry most of the information about ADHD: the
$\delta$ mode alone reaches \R{ModeDeltaSVMBal}\% with the SVM, the $\theta$ mode
\R{ModeThetaSVMBal}\%, while the $\alpha$ (\R{ModeAlphaSVMBal}\%) and $\beta$
(\R{ModeBetaSVMBal}\%) modes are weaker, in line with the many DS-2 windows in which SMVMD finds
no mode in these bands (Table~\ref{tab:commonmodes}). Second, the SMVMD reconstruction gives the
same result as the undecomposed window (\R{ModeReconSVMBal}\%), because the reconstruction
contains almost all of the signal and therefore has nearly the same covariance. Third, fixed
Butterworth bands perform at least as well as the SMVMD common modes overall: the SMVMD
$\delta$ mode is better than the Butterworth $\delta$ band, but in the $\alpha$ and $\beta$
bands, where SMVMD often finds no mode, the filters are better, and with all four components
together the filter bank is slightly ahead. \emph{The adaptive decomposition therefore makes the frequency-specific covariance
representation possible, but it does not by itself add accuracy}; this agrees with the
preliminary finding that features of the raw signal do as well as features of the SMVMD modes
(Chapter~5). Its value in the pipeline lies in a data-driven, interpretable split of the signal
into the oscillations actually present in each window, at an accuracy only one to two points
below the best fixed representations.

\section{Effect of window-level splitting on the proposed pipeline}
To show how the pipeline would be reported under the window-level protocol common in the
literature, the same configuration (all four modes, LightGBM of depth 4) was evaluated with
10-fold cross-validation over windows (Table~\ref{tab:gapmain}). On DS-2 the window-level
protocol reports \R{GapDStwoLeaky}\% against \R{GapDStwoSubj}\% on unseen children, an inflation
of \R{GapDStwoPoints} percentage points. The figures in the 95--100\% range in the literature
(Table~\ref{tab:lit}) are therefore not comparable with subject-wise results, and the proposed
pipeline would itself appear above 90\% if it were evaluated in the same way. On DS-1 the
inflation is small (\R{GapDSoneRestPoints} points) only because the subject-wise result is
already close to 100\%.

\begin{table}[H]
\caption{The proposed pipeline (fixed configuration) under window-wise and subject-wise
cross-validation: balanced window accuracy (\%).}
\label{tab:gapmain}\centering
\gentable{tab_gap_main}\\[2pt]
{\footnotesize \dstwonote}
\end{table}

\section{What the pipeline relies on}
Fig.~\ref{fig:importance} shows the LightGBM importance (total gain, averaged over the 10 outer
DS-2 models) summed per common mode and per electrode, and Table~\ref{tab:topfeat} the ten most
important tangent-space features.

\begin{figure}[H]\centering
\genfig[0.95\textwidth]{importance_main}
\caption{DS-2: share of LightGBM gain per common mode (left) and per electrode (right; the gain
of a channel pair is divided equally between its two electrodes).}
\label{fig:importance}
\end{figure}

\begin{table}[H]
\caption{DS-2: the ten most important tangent-space features of the proposed pipeline. A
``variance'' feature is a diagonal element (power of one channel in the mode), a ``covariance''
feature an off-diagonal element (co-variation of two channels).}
\label{tab:topfeat}\centering
\gentable{tab_top_features}
\end{table}

For ADHD the $\delta$ mode receives \R{ImpDStwoDelta}\% of the importance and the $\theta$ mode
\R{ImpDStwoTheta}\%, against \R{ImpDStwoAlpha}\% and \R{ImpDStwoBeta}\% for the $\alpha$ and $\beta$
modes, in agreement with the single-mode results. Off-diagonal (covariance) features account for
\R{ImpDStwoCovShare}\% of the total importance and diagonal (power) features for
\R{ImpDStwoVarShare}\%. Since the off-diagonal elements are 90\% of the features (171 of 190 per
mode), a single power feature is on average more important than a single channel pair; but most
of the information the classifier uses comes from the co-variation of channels, which per-channel
features cannot capture, consistent with the advantage of connectivity and Riemannian features in
the preliminary study. The most important electrodes are \R{ImpDStwoTopChannels}. Slow activity at frontal electrodes such as Fp2
is also where residual eye-movement artefacts appear, and stimulant medication, which only the
ADHD group received, changes the EEG in the theta band among others~\cite{clarke2002}; the
importance pattern is therefore compatible with both a neural and a non-neural explanation and
should not be read as a biomarker.

\section{Comparison with published results}
Table~\ref{tab:litcomp} places the proposed pipeline beside the published subject-wise results
on the same ADHD dataset (Chapter~2).

\begin{table}[H]
\caption{ADHD dataset: subject-wise results. Published values as reported (accuracy); this work:
balanced accuracy and plain window accuracy.}
\label{tab:litcomp}\centering\small
\begin{adjustbox}{max width=\textwidth}
\begin{tabular}{p{5.6cm}p{3.6cm}p{2.0cm}r}
\hline
Method & Evaluation & Model type & Accuracy (\%) \\
\hline
SVM on spectral features~\cite{geethapriya2026} & subject-wise 10-fold & shallow & 65.6 \\
LSTM~\cite{geethapriya2026} & subject-wise 10-fold & deep (recurrent) & 61.6 \\
CNN--BiLSTM~\cite{geethapriya2026} & subject-wise 10-fold & deep & 65.7 \\
CNN~\cite{geethapriya2026} & subject-wise 10-fold & deep & 72.7 \\
\hline
Starting method, this work (Step~1) & subject-wise 10-fold & shallow & \R{StepADStwoBal} (bal.) \\
\textbf{Proposed pipeline (this work)} & subject-wise 10-fold & \textbf{shallow, CPU} & \textbf{\R{MainDStwoBal}} (bal.) \\
 & & & \R{MainDStwoWinAccPlain} (acc.) \\
\hline
\end{tabular}
\end{adjustbox}\\[2pt]
{\footnotesize \dstwonote}
\end{table}

The proposed pipeline scores higher than the recurrent and hybrid deep networks and the spectral
SVM evaluated subject-wise on the same data, and it is comparable with the best published
subject-wise result, a convolutional network trained end to end, from which it differs by about
one point. The comparison is approximate, since the window length
(5~s with overlap against 4~s without), the folds and the measure (balanced against plain
accuracy) differ. The advantages of the proposed pipeline are of a different kind: it needs no
GPU and no deep network, processes a window in about \R{RtTotal}~ms on one CPU core, trains in
about \R{RtTrain}~s, applies unchanged to a second disorder, and its decisions can be traced to
frequency modes and electrode pairs.

\section{Discussion}
\textbf{Medication and other confounds (DS-2).} The children with ADHD were taking
methylphenidate and the controls were not. Stimulants change the EEG, mainly in the theta and
beta bands~\cite{clarke2002}, and the $\theta$ mode is the second most important mode of the
pipeline. In addition, recordings of children with ADHD are longer (\R{DStwoLenAdhd} against
\R{DStwoLenCtrl} windows on average) because the task lasted until the child answered, and
recording length alone predicts the group with an AUC of \R{DStwoLenAUC}; anything in the signal
that reflects how long or how hard the task was for a child is therefore also a cue. The result
is the detection of medicated ADHD against unmedicated controls, and could partly reflect the
medication or task performance.

\textbf{Uncertainty and selection.} The confidence interval of the main DS-2 result,
\R{MainDStwoBalCI}\%, is about ten points wide, so differences of one or two points between
pipelines are not meaningful. The pipeline's configuration was fixed before its final result was
seen, and its mode set was chosen in the training folds; the per-mode and per-classifier rows of
the ablations are analyses, and the best of them is slightly optimistic. The undecomposed
Riemannian pipeline with LightGBM (\R{AblRawLgbmBal}\%) is numerically the highest single result
of the study, but it is not distinguishable from the proposed pipeline.

\textbf{Diagnostic value.} About three in four unseen children are classified correctly. This is
a substantial and honest improvement over the starting method and comparable with the best
published subject-wise result, but it is far from the accuracy needed for diagnosis, and the
pipeline is a research result, not a clinical tool.
